\appendix
\section{List of Supported \Slice{} Distributions}\label{appen:supported-distributions}
\begin{figure}[!t]
    \centering
    \caption{Distributions supported by \Slice{}, along with their probability density functions.}
    \label{fig:densities}
    \begin{tabular}{l l @{\qquad} l l}
    \toprule
    \textbf{Distribution} & \textbf{Density} & \textbf{Distribution} & \textbf{Density} \\
    \midrule
    Exponential & $\lambda e^{-\lambda x}$ 
        & Uniform & $\frac{1}{b - a}$ \\
    Laplace & $\frac{1}{2b} e^{-\frac{|x - \mu|}{b}}$ 
        & Gaussian & $\frac{1}{\sqrt{2\pi}\sigma} e^{-\frac{(x - \mu)^2}{2\sigma^2}}$ \\
    Cauchy & $\frac{1}{\pi \gamma \left[1 + \left(\frac{x - x_0}{\gamma}\right)^2\right]}$ 
        & Beta & $\frac{x^{\alpha - 1}(1 - x)^{\beta - 1}}{B(\alpha, \beta)}$ \\
    Student-T & $c \left(1 + \frac{x^2}{\nu} \right)^{-\frac{\nu + 1}{2}}$ 
        & LogNormal & $\frac{1}{x\sigma \sqrt{2\pi}} e^{-\frac{(\ln x - \mu)^2}{2\sigma^2}}$ \\
    Chi-squared & $\frac{1}{2^{k/2} \Gamma(k/2)} x^{k/2 - 1} e^{-x/2}$ 
        & Gamma & $\frac{\beta^\alpha}{\Gamma(\alpha)} x^{\alpha - 1} e^{-\beta x}$ \\
    Logistic & $\frac{e^{-(x - \mu)/s}}{s(1 + e^{-(x - \mu)/s})^2}$ 
        & Pareto & $\frac{\alpha x_m^\alpha}{x^{\alpha + 1}}$ \\
    Rayleigh & $\frac{x}{\sigma^2} e^{-x^2 / (2\sigma^2)}$ 
        & Weibull & $\frac{k}{\lambda} \left( \frac{x}{\lambda} \right)^{k-1} e^{-(x/\lambda)^k}$ \\
    & & Gumbel-1 & $\frac{1}{\beta} e^{-(z + e^{-z})}$,\quad $z = \frac{x - \mu}{\beta}$ \\
    & & Gumbel-2 & $\frac{\alpha \beta x^{-\alpha - 1} e^{-(x/\beta)^{-\alpha}}}{\beta}$ \\
    & & Exponentiated Power & $k \lambda x^{k - 1} e^{-(x/\lambda)^k}$ \\
    \bottomrule
    \end{tabular}
    \end{figure}
    

\section{Typing Rules for Full \Slice{} Language}\label{appen:typing-rules}
\begin{figure}
\begin{mathpar}
    \inferrule[\textsc{Var}]
    {\ }
    {\Gamma, x: \tau \vdash x : \tau}

    \inferrule[\textsc{Let}]
    {\Gamma \vdash e_1 : \tau_1 \\
     \Gamma, x: \tau_1 \vdash e_2 : \tau_2}
    {\Gamma \vdash \letkw \; x = e_1 \; \inkw \; e_2 : \tau_2}

    \inferrule[\textsc{If}]
    {\Gamma \vdash e_1 : \bool \\
     \Gamma \vdash e_2 : \tau \\
     \Gamma \vdash e_3 : \tau}
    {\Gamma \vdash \ifkw \; e_1 \; \thenkw \; e_2 \; \elsekw \; e_3 : \tau}

    \inferrule[\textsc{Discrete}]
    {n \geq 1 \\
     \forall i \in \{0,\ldots,n-1\}.\; p_i \geq 0 \\
     \sum_{i=0}^{n-1} p_i = 1}
    {\Gamma \vdash \discrete(p_0, \ldots, p_{n-1}) : \fin{n}}

    \inferrule[\textsc{LessEq}]
    {\Gamma \vdash e : \intty}
    {\Gamma \vdash e \leq i : \bool}

    \inferrule[\textsc{Pair}]
    {\Gamma \vdash e_1 : \tau_1 \\
     \Gamma \vdash e_2 : \tau_2}
    {\Gamma \vdash (e_1, e_2) : \tau_1 * \tau_2}

    \inferrule[\textsc{Fst}]
    {\Gamma \vdash e : \tau_1 * \tau_2}
    {\Gamma \vdash \fstkw \; e : \tau_1}

    \inferrule[\textsc{Snd}]
    {\Gamma \vdash e : \tau_1 * \tau_2}
    {\Gamma \vdash \sndkw \; e : \tau_2}

    \inferrule[\textsc{Fun}]
    {\Gamma, x: \tau_1 \vdash e : \tau_2}
    {\Gamma \vdash \funkw \; x \; \rightarrow \; e : \tau_1 \rightarrow \tau_2}

    \inferrule[\textsc{App}]
    {\Gamma \vdash e_1 : \tau_1 \rightarrow \tau_2 \\
     \Gamma \vdash e_2 : \tau_1}
    {\Gamma \vdash e_1 \; e_2 : \tau_2}

    \inferrule[\textsc{True}]
    {\ }
    {\Gamma \vdash \text{true} : \bool}

    \inferrule[\textsc{False}]
    {\ }
    {\Gamma \vdash \text{false} : \bool}

    \inferrule[\textsc{And}]
    {\Gamma \vdash e_1 : \bool \\
     \Gamma \vdash e_2 : \bool}
    {\Gamma \vdash e_1 \logand e_2 : \bool}

    \inferrule[\textsc{Or}]
    {\Gamma \vdash e_1 : \bool \\
     \Gamma \vdash e_2 : \bool}
    {\Gamma \vdash e_1 \logor e_2 : \bool}

    \inferrule[\textsc{Not}]
    {\Gamma \vdash e : \bool}
    {\Gamma \vdash \text{not}\; e : \bool}

    \inferrule[\textsc{Unit}]
    {\ }
    {\Gamma \vdash () : \unit}

    \inferrule[\textsc{FinConst}]
    {0 \leq k < n}
    {\Gamma \vdash \finconst{k}{n} : \fin{n}}

    \inferrule[\textsc{FinLess}]
    {\Gamma \vdash e_1 : \fin{n} \\
     \Gamma \vdash e_2 : \fin{n}}
    {\Gamma \vdash e_1 \finlt{n} e_2 : \bool}

    \inferrule[\textsc{FinLessEq}]
    {\Gamma \vdash e_1 : \fin{n} \\
     \Gamma \vdash e_2 : \fin{n}}
    {\Gamma \vdash e_1 \finleq{n} e_2 : \bool}

    \inferrule[\textsc{Observe}]
    {\Gamma \vdash e : \bool}
    {\Gamma \vdash \observekw\; e : \unit}

    \inferrule[\textsc{Seq}]
    {\Gamma \vdash e_1 : \tau_1 \\
     \Gamma \vdash e_2 : \tau_2}
    {\Gamma \vdash e_1; e_2 : \tau_2}

    \inferrule[\textsc{Fix}]
    {\Gamma, f: \tau_1 \rightarrow \tau_2, x: \tau_1 \vdash e : \tau_2}
    {\Gamma \vdash \text{fix}\; f\; x := e : \tau_1 \rightarrow \tau_2}

    \inferrule[\textsc{Nil}]
    {\ }
    {\Gamma \vdash \text{nil} : \listty{\tau}}

    \inferrule[\textsc{Cons}]
    {\Gamma \vdash e_1 : \tau \\
     \Gamma \vdash e_2 : \listty{\tau}}
    {\Gamma \vdash e_1 :: e_2 : \listty{\tau}}

    \inferrule[\textsc{Match}]
    {\Gamma \vdash e : \listty{\tau_1} \\
     \Gamma \vdash e_1 : \tau_2 \\
     \Gamma, h: \tau_1, t: \listty{\tau_1} \vdash e_2 : \tau_2}
    {\Gamma \vdash \text{match}\; e\; \text{with}\; \text{nil} \rightarrow e_1 \mid h :: t \rightarrow e_2\; \text{end} : \tau_2}

    % \inferrule[\textsc{Ref}]
    % {\Gamma \vdash e : \tau}
    % {\Gamma \vdash \text{ref}\; e : \text{ref}(\tau)}

    % \inferrule[\textsc{Deref}]
    % {\Gamma \vdash e : \text{ref}(\tau)}
    % {\Gamma \vdash !e : \tau}

    % \inferrule[\textsc{Assign}]
    % {\Gamma \vdash e_1 : \text{ref}(\tau) \\
    %  \Gamma \vdash e_2 : \tau}
    % {\Gamma \vdash e_1 := e_2 : \unit}

    \inferrule[\textsc{Diverge}]
    {\ }
    {\Gamma \vdash \diverge : \tau}
\end{mathpar}
\caption{Typing rules for general language constructs}
\label{fig:typing-general-all}
\end{figure}


% \section{Proofs for Distinguishability}
% \subsection{Proof of Lemma~\ref{thm:distinguishes-singleton}}
% \begin{proof}
% 	Follows directly from Definition~\ref{def:distinguishes}. Since $V$ is a singleton set, and there can be no two elements that lie in the same interval, it must be the case that $\intervals{B} = \{(-\infty, +\infty)\}$, which means that $B = \emptyset$.
% \end{proof}

% \subsection{Proof of Lemma~\ref{thm:distinguishes-top}}
% \begin{proof}
% 	Suppose by contradiction that $B \neq \top$ but $B$ distinguishes $V = \top$. Since $B \neq \top$, the set $\intervals{B}$ partitions $\mathbb{R}$ into finitely many intervals. However, since $V = \top$ represents the set of all real numbers, and there are uncountably many real values, it is impossible for every interval in $\intervals{B}$ to contain at most one element from $V$. In particular, at least one interval must contain more than one value from $V$, violating the condition in Definition~\ref{def:distinguishes}. Therefore, $B = \top$.
% \end{proof}

\section{Proofs for Discretization}
\subsection{Proof of Lemma~\ref{lem:type-transformations}}
\begin{proof}\label{proof:type-transformations}
    We proceed by induction on the structure of the typing derivation $\Gamma \vdash e : \tau$. That is, we match each typing rule and show that it remains valid after applying $\discretize{\cdot}$ to the context, expression, and type.
    
    \textbf{Case 1:} $e = x$
    \begin{itemize}
        \item We have that $\discretize{x : \tau} = x$, and since $\Gamma(x) = \tau$, we have that $\discretize{\Gamma}(x) = \discretize{\tau} = \tau$ for $x \in \text{dom}(\Gamma)$. Thus, $\discretize{\Gamma} \vdash x : \discretize{\tau}$.
    \end{itemize}

    \textbf{Case 2:} $e = (c : \float[B;V])$
    \begin{itemize}
        \item If $B = \top$, then $\discretize{c : \float[B;V]} = c$, and $\discretize{\float[B;V]} = \float$. As per the float typing rules in Figure~\ref{fig:typing-float}, $c$ is a constant of type $\float$.
        \item If $B \neq \top$, then $\discretize{e : \float[B;V]} = \finconst{k}{n+1}$ for some $k$ such that $c \in \intervals{B}_k$, and $\discretize{\float[B;V]} = \fin{n+1}$. As per the typing rules in Figure~\ref{fig:typing-general}, $\finconst{k}{n+1}$ is a constant of type $\fin{n+1}$.
    \end{itemize}

    \textbf{Case 3:} $e = (\letkw \; x = e_1 : \tau_1 \; \inkw \; e_2 : \tau_2)$
    \begin{itemize}
        \item The \textsc{Let} typing rule is shown in Figure~\ref{fig:typing-general}. By the inductive hypothesis, we assume $\discretize{\Gamma} \vdash \discretize{e_1 : \tau_1} : \discretize{\tau_1}$ and $\discretize{\Gamma, x : \tau_1} \vdash \discretize{e_2 : \tau_2} : \discretize{\tau_2}$. Then, we can conclude:
        \[ \discretize{\Gamma} \vdash \letkw \; x = \discretize{e_1 : \tau_1} \; \inkw \; \discretize{e_2 : \tau_2} : \discretize{\tau_2}. \]
    \end{itemize}

    \textbf{Case 4:} $e = (e_1 : \float[B_1; V_1] < e_2 : \float[B_2; V_2])$
    \begin{itemize}
        \item The \textsc{Less} typing rule is shown in Figure~\ref{fig:typing-float}. By the inductive hypothesis, we assume $\discretize{\Gamma} \vdash \discretize{e_1 : \tau_1} : \discretize{\tau_1}$ and $\discretize{\Gamma} \vdash \discretize{e_2 : \tau_2} : \discretize{\tau_2}$, where $\tau_1$ and $\tau_2$ are $\float[B_1;V_1]$ and $\float[B_2;V_2]$, respectively. Since both discretized expressions are either floats or finite-typed values, the comparison remains well-typed under standard typing rules:
        \[
        \discretize{\Gamma} \vdash \discretize{e_1 : \float[B_1; V_1]} < \discretize{e_2 : \float[B_2; V_2]} : \discretize{\bool}.
        \]
    \end{itemize}

    \textbf{Case 5:} $e = (\ifkw\; e_1 : \bool \; \thenkw\; e_2 : \tau \; \elsekw\; e_3 : \tau)$

    \begin{itemize}
        \item The \textsc{If} typing rule is shown in Figure~\ref{fig:typing-general}. By the inductive hypothesis, we assume $\discretize{\Gamma} \vdash \discretize{e_1 : \bool} : \discretize{\bool}$, $\discretize{\Gamma} \vdash \discretize{e_2 : \tau} : \discretize{\tau}$, and $ \discretize{\Gamma} \vdash \discretize{e_3 : \tau} : \discretize{\tau}$. Then, we can conclude:
        \[
            \discretize{\Gamma} \vdash \ifkw\; \discretize{e_1 : \bool} \; \thenkw\; \discretize{e_2 : \tau} \; \elsekw\; \discretize{e_3 : \tau} : \discretize{\tau}
        \]
    \end{itemize}

    \textbf{Case 6:} $e = (\uniform(e_1, e_2) : \float[B;V])$

    \begin{itemize}
        \item If $B = \top$, then $\discretize{\float[B;V]} = \float$, and thus:
        \[
            \discretize{\uniform(e_1, e_2) : \float[B;V]} = \uniform(\discretize{e_1 : \tau_1}, \discretize{e_2 : \tau_2}) : \float
        \]
        and by the inductive hypothesis, both $e_1$ and $e_2$ with their respective types $\tau_1 = \float[B_1; V_1]$ and $\tau_2 = \float[B_2; V_2]$ are well-typed, so the result is well-typed under the float rule.

        \item If $B \neq \top$ and $e_1, e_2 \in \mathbb{R}$, then $\discretize{\float[B;V]} = \fin{n+1}$, and thus:
        \[
            \discretize{\uniform(e_1, e_2) : \float[B;V]} = \discrete(p_0, \ldots, p_n) : \fin{n+1}
        \]
        and each $p_k = \text{CDF}(b_{k+1}) - \text{CDF}(b_k)$ satisfies $p_k \in [0,1]$ and $\sum p_k = 1$, so this is well-typed as a finite distribution.

        \item If $B \neq \top$, and $e_1$ or $e_2$ is variable, then the algorithm generates conditionals to evaluate the possible values from $V_1$ and $V_2$ using fresh variables $v_1'$ and $v_2'$, and dispatches to an appropriate discretization of $\uniform$ with constant parameters in each branch. Each branch is well-typed by the inductive hypothesis, and the full expression is well-typed via the typing rules for conditionals and finite types.
    \end{itemize}

    % \textbf{Case 7:} Compound types (functions, pairs, lists, references)

    % \begin{itemize}
    %     \item These follow directly by applying $\discretize{\cdot}$ recursively on the subcomponents of the type and expression. By the inductive hypothesis on each subterm, the transformed expression is well-typed under the transformed type. For instance:
    %     \begin{itemize}
    %         \item If $\tau = A \to B$, then $\discretize{\tau} = \discretize{A} \to \discretize{B}$ and function abstraction preserves typing.
    %         \item If $\tau = A \times B$, then $\discretize{e} = (\discretize{e_1}, \discretize{e_2})$ is typed under $\discretize{A} \times \discretize{B}$.
    %         \item Likewise for lists $\listty{A}$ and references $\texttt{ref}(A)$.
    %     \end{itemize}
    % \end{itemize}

    Hence, each case of the typing derivation is preserved under $\discretize{\cdot}$, which acts on expressions $e : \tau$, types $\tau$, and the context $\Gamma$. Therefore,
    \[
    \discretize{\Gamma} \vdash \discretize{e : \tau} : \discretize{\tau}.
    \]

\end{proof}

\section{Soundness Proof}\label{sec:soundness}

This section establishes the soundness of our discretization transformation. We proceed in three steps: first, we define a nondeterministic small-step semantics for the base language; second, we extend this to a probabilistic semantics that tracks distributions; and third, we prove that discretization preserves the probabilistic semantics.

\subsection{Nondeterministic Small-Step Semantics}

We begin by defining a small-step operational semantics for \Slice{} that treats probabilistic sampling nondeterministically. This semantics captures the possible values that expressions can take without tracking their probabilities.

\subsubsection{Values}

First, we define the set of values that expressions can evaluate to:
\begin{align*}
v ::= &\; c                                  & \text{float constant} \\
    | &\; \text{true} \mid \text{false}      & \text{boolean value} \\
    | &\; \finconst{k}{n}                    & \text{finite type value} \\
    | &\; i                                  & \text{integer value} \\
    | &\; ()                                 & \text{unit value} \\
    | &\; (v_1, v_2)                         & \text{pair value} \\
    | &\; \funkw \; x \; \rightarrow \; e    & \text{function closure} \\
    | &\; \text{nil}                         & \text{empty list} \\
    | &\; v_1 :: v_2                         & \text{list cons value} \\
    | &\; \ell                               & \text{location (for references)}
\end{align*}

\subsubsection{Evaluation Contexts}

We use evaluation contexts to specify the order of evaluation:
\begin{align*}
E ::= &\; [\cdot] \\
    | &\; \letkw \; x = E \; \inkw \; e \\
    | &\; E < e \mid v < E \\
    | &\; E \leq e \mid v \leq E \\
    | &\; E \finlt{n} e \mid v \finlt{n} E \\
    | &\; E \finleq{n} e \mid v \finleq{n} E \\
    | &\; E \logand e \mid v \logand E \\
    | &\; E \logor e \mid v \logor E \\
    | &\; \text{not}\; E \\
    | &\; \ifkw \; E \; \thenkw \; e_1 \; \elsekw \; e_2 \\
    | &\; (E, e) \mid (v, E) \\
    | &\; \fstkw \; E \mid \sndkw \; E \\
    | &\; E \; e \mid v \; E \\
    | &\; \observekw\; E \\
    | &\; E; e \\
    | &\; E :: e \mid v :: E \\
    | &\; \text{match}\; E \; \text{with}\; \ldots \\
    | &\; \text{ref}\; E \\
    | &\; !E \\
    | &\; E := e \mid v := E
\end{align*}

\subsubsection{Small-Step Relation}

We define the small-step relation $e \rightarrow \mathcal{P}(e')$ that maps an expression to a set of possible next expressions. For deterministic constructs, this set is a singleton; for probabilistic sampling, it contains all possible values in the distribution's support.

\paragraph{Basic Rules}

\begin{mathpar}
\inferrule[\textsc{Let}]
{\ }
{\letkw \; x = v \; \inkw \; e \rightarrow \{e[v/x]\}}

\inferrule[\textsc{IfTrue}]
{\ }
{\ifkw \; \text{true} \; \thenkw \; e_1 \; \elsekw \; e_2 \rightarrow \{e_1\}}

\inferrule[\textsc{IfFalse}]
{\ }
{\ifkw \; \text{false} \; \thenkw \; e_1 \; \elsekw \; e_2 \rightarrow \{e_2\}}

\inferrule[\textsc{Fst}]
{\ }
{\fstkw \; (v_1, v_2) \rightarrow \{v_1\}}

\inferrule[\textsc{Snd}]
{\ }
{\sndkw \; (v_1, v_2) \rightarrow \{v_2\}}

\inferrule[\textsc{App}]
{\ }
{(\funkw \; x \; \rightarrow \; e) \; v \rightarrow \{e[v/x]\}}
\end{mathpar}

\paragraph{Comparison Rules}

\begin{mathpar}
\inferrule[\textsc{LessTrue}]
{c_1 < c_2}
{c_1 < c_2 \rightarrow \{\text{true}\}}

\inferrule[\textsc{LessFalse}]
{c_1 \geq c_2}
{c_1 < c_2 \rightarrow \{\text{false}\}}

\inferrule[\textsc{LeqTrue}]
{i_1 \leq i_2}
{i_1 \leq i_2 \rightarrow \{\text{true}\}}

\inferrule[\textsc{LeqFalse}]
{i_1 > i_2}
{i_1 \leq i_2 \rightarrow \{\text{false}\}}
\end{mathpar}

\paragraph{Boolean Operations}

\begin{mathpar}
\inferrule[\textsc{AndTrue}]
{\ }
{\text{true} \logand e \rightarrow \{e\}}

\inferrule[\textsc{AndFalse}]
{\ }
{\text{false} \logand e \rightarrow \{\text{false}\}}

\inferrule[\textsc{OrTrue}]
{\ }
{\text{true} \logor e \rightarrow \{\text{true}\}}

\inferrule[\textsc{OrFalse}]
{\ }
{\text{false} \logor e \rightarrow \{e\}}

\inferrule[\textsc{NotTrue}]
{\ }
{\text{not}\; \text{true} \rightarrow \{\text{false}\}}

\inferrule[\textsc{NotFalse}]
{\ }
{\text{not}\; \text{false} \rightarrow \{\text{true}\}}
\end{mathpar}

\paragraph{List Operations}

\begin{mathpar}
\inferrule[\textsc{MatchNil}]
{\ }
{\text{match}\; \text{nil} \; \text{with}\; \text{nil} \rightarrow e_1 \mid h :: t \rightarrow e_2 \; \text{end} \rightarrow \{e_1\}}

\inferrule[\textsc{MatchCons}]
{\ }
{\text{match}\; v_1 :: v_2 \; \text{with}\; \text{nil} \rightarrow e_1 \mid h :: t \rightarrow e_2 \; \text{end} \rightarrow \{e_2[v_1/h, v_2/t]\}}
\end{mathpar}

\paragraph{Probabilistic Sampling (Nondeterministic)}

For continuous distributions, the small-step relation returns the set of all possible values in the distribution's support:

\begin{mathpar}
\inferrule[\textsc{Uniform}]
{v_1, v_2 \text{ are float values} \\ v_1 < v_2}
{\uniform(v_1, v_2) \rightarrow \{c \mid v_1 \leq c < v_2\}}

\inferrule[\textsc{Gaussian}]
{\mu, \sigma \text{ are float values} \\ \sigma > 0}
{\gaussian(\mu, \sigma) \rightarrow \{c \mid c \in \mathbb{R}\}}

\inferrule[\textsc{Exponential}]
{\lambda \text{ is a float value} \\ \lambda > 0}
{\exponential(\lambda) \rightarrow \{c \mid c \geq 0\}}

\inferrule[\textsc{Beta}]
{\alpha, \beta \text{ are float values} \\ \alpha > 0 \\ \beta > 0}
{\betafn(\alpha, \beta) \rightarrow \{c \mid 0 \leq c \leq 1\}}

\inferrule[\textsc{Discrete}]
{p_0, \ldots, p_n \text{ are probabilities}}
{\discrete(p_0, \ldots, p_n) \rightarrow \{0, 1, \ldots, n\}}
\end{mathpar}

\paragraph{Observations}

\begin{mathpar}
\inferrule[\textsc{ObserveTrue}]
{\ }
{\observekw\; \text{true} \rightarrow \{()\}}

\inferrule[\textsc{ObserveFalse}]
{\ }
{\observekw\; \text{false} \rightarrow \emptyset}
\end{mathpar}

Note that observing a false condition results in an empty set, representing program termination without producing a value.

\paragraph{References and State}

For references, we extend the semantics with a store $\sigma$ that maps locations to values. The relation becomes $\langle e, \sigma \rangle \rightarrow \mathcal{P}(\langle e', \sigma' \rangle)$:

\begin{mathpar}
\inferrule[\textsc{Ref}]
{\ell \text{ is a fresh location}}
{\langle \text{ref}\; v, \sigma \rangle \rightarrow \{\langle \ell, \sigma[\ell \mapsto v] \rangle\}}

\inferrule[\textsc{Deref}]
{\sigma(\ell) = v}
{\langle !\ell, \sigma \rangle \rightarrow \{\langle v, \sigma \rangle\}}

\inferrule[\textsc{Assign}]
{\ }
{\langle \ell := v, \sigma \rangle \rightarrow \{\langle (), \sigma[\ell \mapsto v] \rangle\}}
\end{mathpar}

\paragraph{Context Rule}

\begin{mathpar}
\inferrule[\textsc{Context}]
{e \rightarrow S}
{E[e] \rightarrow \{E[e'] \mid e' \in S\}}
\end{mathpar}

This nondeterministic semantics captures all possible execution paths of a probabilistic program without tracking the probability of each path. In the next subsection, we will extend this to a probabilistic semantics that properly tracks distributions.

\subsection{Probabilistic Small-Step Semantics}

We now define a probabilistic semantics that tracks the distribution over expressions. Define
%
\[
	\extexprs = \exprs \cup  \{ \error \} \cup \{ \obsfail \}~,
\]
%
and call them \emph{extended expressions}.
%
%where $\error$ indicates a type error and $\obsfail$ denotes observation failure.
The semantic function $\sem{\cdot} \colon \exprs \rightarrow \distexprs$ maps expressions to distributions over expressions, where $\dist{X}$ denotes the space of probability distributions over $X$.

We assume the existence of the Giry monad structure on $\mathcal{D}$, which provides:
\begin{itemize}
    \item $\monunit \colon X \rightarrow \dist{X}$ (also written as $\delta_x$ for the Dirac distribution at $x$)
    \item $\monbind \colon \dist{X} \times (X \to \dist{Y}) \rightarrow \dist{Y}$ (monadic bind)
\end{itemize}

\subsubsection{Deterministic Constructs}

For deterministic constructs, the semantics returns a Dirac distribution:

\paragraph{Values and Basic Operations}
\begin{align*}
\sem{\val} &= \delta_\val \quad \text{(for any value } \val \text{)} \\
\sem{\letkw \; \var = \val \; \inkw \; \expr} &= \delta_{\expr\subst{\val}{\var}} \\
%
%
\sem{ \fstkw \; \val} & = 
\begin{cases}
	\delta_{\val_1}  & \text{if $\val = (\val_1,\val_2)$ for some values $\val_1,\val_2$} \\
	%
	\delta_{\error}  & \text{otherwise}
\end{cases}
%
\\
%
%
\sem{ \sndkw \; \val} & = 
\begin{cases}
	\delta_{\val_2}  & \text{if $\val = (\val_1,\val_2)$ for some values $\val_1,\val_2$} \\
	%
	\delta_{\error}  & \text{otherwise}
\end{cases}
%
\\
%
\sem{ (\funkw \; \var \; \rightarrow \; \expr) \; \val } &= \delta_{\expr\subst{\val}{\var}}\\
\sem{\error} = \delta_\error \\
\sem{\obsfail} = \delta_\obsfail
\end{align*}
%

\paragraph{Conditionals}
\begin{align*}
\sem{ \ifkw \; \val \; \thenkw \; e_1 \; \elsekw \; e_2 } & = 
\begin{cases}
	\sem{\expr_1}  & \text{if $\val = \true$} \\
	%
	\sem{\expr_2}  & \text{if $\val = \false$} \\
	%
	\delta_\error & \text{otherwise}
\end{cases}
\end{align*}
%

\paragraph{Comparisons}
\begin{align*}
\sem{ v_1 < v_2 } &= \begin{cases}
    \delta_{\text{true}} & \text{if $v_1$ and $v_2$ are float constants and~} v_1 < v_2 \\
    \delta_{\text{false}} & \text{if $v_1$ and $v_2$ are float constants and~} v_1 \geq v_2\\
    \delta_\error & \text{otherwise}
\end{cases} \\
\sem{ v_1 \leq v_2 } &=\begin{cases}
	\delta_{\text{true}} & \text{if $v_1$ and $v_2$ are integer values and~} v_1 < v_2 \\
	\delta_{\text{false}} & \text{if $v_1$ and $v_2$ are integer values and~} v_1 \geq v_2\\
	\delta_\error & \text{otherwise}
\end{cases}
\end{align*}
%

\paragraph{Boolean Operations}
\begin{align*}
%\sem{ \text{true} \logand e } &=  \sem{ e } \\
%\sem{ \text{false} \logand e } &= \delta_{\text{false}} \\
%
\sem{ v \logand e } & = \begin{cases}
	\sem{ e } & \text{if $v=\true$} \\
	\delta_{\text{false}} & \text{if $v=\false$} \\
	\delta_\error & \text{otherwise}
\end{cases} \\ 
%
%=
%\sem{ \text{true} \logor e } &= \delta_{\text{true}} \\
%\sem{ \text{false} \logor e } &= \sem{ e } \\
%
\sem{ v \logor e } & = \begin{cases}
	\delta_\true & \text{if $v=\true$} \\
	\sem{ e } & \text{if $v=\false$} \\
	\delta_\error & \text{otherwise}
\end{cases} \\ 
%
%
\sem{ \text{not}\; v } &= 
 \begin{cases}
	\delta_\false & \text{if $v=\true$} \\
	\delta_\true & \text{if $v=\false$} \\
	\delta_\error & \text{otherwise}
\end{cases} \\ 
 \\
%
%\sem{ \text{not}\; \text{true} } &= \delta_{\text{false}} \\
%\sem{ \text{not}\; \text{false} } &= \delta_{\text{true}}
\end{align*}
%


\subsubsection{Compositional Rules}

For compound expressions, we use monadic bind to compose distributions:

\paragraph{Binary Operations}
\begin{align*}
\sem{ e_1 < e_2 } &= \sem{ e_1 } \gg\!= \lambda v_1. \sem{ v_1 < e_2 } \\
\sem{ v_1 < e_2 } &= \sem{ e_2 } \gg\!= \lambda v_2. \sem{ v_1 < v_2 }
\end{align*}

Similar rules apply for $\leq$, $\logand$, $\logor$, and other binary operations.

\paragraph{Let Binding}
\begin{align*}
\sem{ \letkw \; x = e_1 \; \inkw \; e_2 } &= \sem{ e_1 } \gg\!= \lambda v. \sem{ \letkw \; x = v \; \inkw \; e_2 }
\end{align*}

\paragraph{Conditionals}
\begin{align*}
\sem{ \ifkw \; e_1 \; \thenkw \; e_2 \; \elsekw \; e_3 } &= \sem{ e_1 } \gg\!= \lambda v. \sem{ \ifkw \; v \; \thenkw \; e_2 \; \elsekw \; e_3 }
\end{align*}

\paragraph{Function Application}
\begin{align*}
\sem{ e_1 \; e_2 } &= \sem{ e_1 } \gg\!= \lambda v_1. \sem{ v_1 \; e_2 } \\
\sem{ v_1 \; e_2 } &= \sem{ e_2 } \gg\!= \lambda v_2. \sem{ v_1 \; v_2 }
\end{align*}

\paragraph{Pairs and Projections}
\begin{align*}
\sem{ (e_1, e_2) } &= \sem{ e_1 } \gg\!= \lambda v_1. \sem{ (v_1, e_2) } \\
\sem{ (v_1, e_2) } &= \sem{ e_2 } \gg\!= \lambda v_2. \delta_{(v_1, v_2)} \\
\sem{ \fstkw \; e } &= \sem{ e } \gg\!= \lambda v. \sem{ \fstkw \; v } \\
\sem{ \sndkw \; e } &= \sem{ e } \gg\!= \lambda v. \sem{ \sndkw \; v }
\end{align*}

\paragraph{Sequencing}
\begin{align*}
\sem{ e_1; e_2 } &= \sem{ e_1 } \gg\!= \lambda v_1. \sem{ e_2 }
\end{align*}

\paragraph{Finite Type Operations}
\kevin{why don't we distinguish values and compound expressions here?}
\begin{align*}
\sem{ \finconst{k}{n} } &= \delta_{k_{\#n}} \\
%
\sem{ e_1 \finlt{n} e_2 } &= \sem{ e_1 } \gg\!= \lambda v_1. \sem{ v_1 \finlt{n} e_2 } \\
%
\sem{ k_{1,\#n} \finlt{n} k_{2,\#n} } &= \begin{cases}
    \delta_{\text{true}} & \text{if } k_1 < k_2 \\
    \delta_{\text{false}} & \text{otherwise}
\end{cases} \\
%
\sem{ v_1 \finlt{n} v_1 } &= \begin{cases}
	\delta_{\text{true}} & \text{if $v_1 = k_{1,\#n}$ and $v_2 = k_{1,\#n}'$ are finite type values and}~k < k' \\
	\delta_{\text{false}} & \text{if $v_1 = k_{1,\#n}$ and $v_2 = k_{1,\#n}'$ are finite type values and}~k \geq k' \\
	\delta_{\error} & \text{otherwise}
\end{cases}
\end{align*}

\subsubsection{Probabilistic Distributions}

For sampling expressions, we map to the corresponding probability distributions. Undefined sampling is interpreted as divergence. 

\begin{align*}
\sem{ \uniform(v_1, v_2) } &= \begin{cases}
\text{Uniform}_{[v_1, v_2)}  & \text{if $v_1$ and $v_2$ are float constants with $v_1 < v_2$} \\
%
\delta_{\uniform(v_1, v_2)} & \text{otherwise}
\end{cases}
\\
\sem{ \gaussian(v_\mu, v_\sigma) } &= 
 \begin{cases}
	\mathcal{N}(v_\mu, v_\sigma^2)  & \text{if $v_\mu$ and $v_\sigma$ are float constants with $v_\sigma > 0$} \\
	%
	\delta_{\gaussian(v_\mu, v_\sigma)} & \text{otherwise}
\end{cases}
\\
\sem{ \exponential(v_\lambda) } &= 
 \begin{cases}
	\text{Exp}(v_\lambda)   & \text{if $v_\lambda$ is a float constant with $v_\lambda > 0$} \\
	%
	\delta_{\exponential(v_\lambda)} & \text{otherwise}
\end{cases}
\\
\sem{ \betafn(v_\alpha, v_\beta) } &= 
 \begin{cases}
	\text{Beta}(v_\alpha, v_\beta)   & \text{if $v_\alpha$ and $v_\beta$ are float constants with $v_\alpha > 0$ and $v_\beta > 0$} \\
	%
	\delta_{\betafn(v_\alpha, v_\beta)} & \text{otherwise}
\end{cases}
\\
\sem{ \discrete(v_0, \ldots, v_n) } &= %\sum_{i=0}^n p_i \cdot \delta_i
\begin{cases}
	\sum_{i=0}^n v_i \cdot \delta_i & \text{if the $v_i$ a are non-negative float constants with $\sum_i v_i = 1$} \\
	%
	\delta_{\discrete(v_0, \ldots, v_n)}  & \text{otherwise}
\end{cases}
\end{align*}
\kevin{double check discrete}

Note that for continuous distributions, we have an implicit coercion from distributions over reals to distributions over constant expressions. That is, if $D$ is a distribution over $\mathbb{R}$, we treat it as a distribution over expressions by mapping each real $r$ to the constant expression $r$.

\subsubsection{Observations}

Observations condition the distribution:

\begin{align*}
\sem{ \observekw\; e } &= \sem{ e } \gg\!= \lambda v. \sem{ \observekw\; v } \\
%\sem{ \observekw\; \text{true} } &= \delta_{()} \\
%\sem{ \observekw\; \text{false} } &= \bot
%
\sem{ \observekw\; v} &=
\begin{cases}
	\delta_{()}  & \text{if $v=\true$} \\
	%
	\delta_{\obsfail} & \text{if $v=\false$} \\
	%
	\delta_\error & \text{otherwise}
\end{cases}
\end{align*}

where $\bot$ represents the zero measure (failed observation).

\subsubsection{DistrCase}

For distrcase expressions (which represent explicit discrete distributions in the implementation):

\begin{align*}
\sem{ \text{distrcase}\{(e_0, p_0), \ldots, (e_n, p_n)\} } &= \sum_{i=0}^n p_i \cdot \sem{ e_i }
\end{align*}

\subsubsection{Example}

Consider the expression:
\begin{lstlisting}
let x = uniform(0.0, 1.0) in
if x < 0.5 then 1 else 2
\end{lstlisting}

The probabilistic semantics gives:
\begin{align*}
&\sem{ \text{let } x = \uniform(0.0, 1.0) \text{ in if } x < 0.5 \text{ then } 1 \text{ else } 2 } \\
&= \sem{ \uniform(0.0, 1.0) } \gg\!= \lambda x. \sem{ \text{if } x < 0.5 \text{ then } 1 \text{ else } 2 } \\
&= \text{Uniform}_{[0,1)} \gg\!= \lambda x. \sem{ \text{if } x < 0.5 \text{ then } 1 \text{ else } 2 } \\
&= \text{Uniform}_{[0,1)} \gg\!= \lambda x. \begin{cases}
    \delta_1 & \text{if } x < 0.5 \\
    \delta_2 & \text{if } x \geq 0.5
\end{cases} \\
&= 0.5 \cdot \delta_1 + 0.5 \cdot \delta_2
\end{align*}

%This probabilistic semantics forms the foundation for proving that our discretization transformation preserves program behavior.
%%
%\begin{theorem}
%	\label{thm:small-step-well-defined}
%	For every $\expr \in \exprs$, $\sem{\expr}$ is a well defined distribution in $\distextexprs$.
%\end{theorem}
%%
%\begin{proof}
%	By induction on $\exprs$.
%\end{proof}
%
%
\kevin{Let's see whether nondet semantics requires similar treatment of undefied sampling.}
\kevin{keep an eye on everything being measurable (will mostly be trivial due to nice construction of measure space)}


\begin{theorem}[Almost-Sure Progress and Preservation]
	\label{thm:small-step-progress-preservation} We have 
	%
	\begin{itemize}
		\item Almost-sure progress: If $\expr : \tau$, then $\sem{\expr}(\{ \error \}) = 0$.
		%
		\item Almost-sure preservation: If $\expr : \tau$, then $\sem{\expr}(\{ \expr' ~|~ \expr' : \tau \}) = 1$.
	\end{itemize}
	
\end{theorem}
%
\begin{proof}
	\kevin{probably} By induction on the derivation tree for $\expr : \tau$.
\end{proof}

\subsection{From Small-Step to Big-Step Semantics}
%
%
%
For every $\expr \in \exprs$ and all $n \in \mathbb{N}$, we define the distribution $\bigsemstep{\expr}{n} \in \distextexprs$ recursively as
%
\[
	\bigsemstep{\expr}{n} = 
	\begin{cases}
		\delta_\expr & \text{if $n = 0$} \\
		%
		\bigsemstep{\expr}{n-1} \monbind  \lambda \expr'. \sem{\expr'} & \text{if $n > 0$}~.
	\end{cases}
\]
%
Let $\vals$ be the measurable set of all values and write $\nonvals = \expr \setminus \vals$, which is also measurable.
%
%
\begin{lemma}[Basic Properties of $\bigsemstep{\cdot}{n}$]
	\label{lem:props:bigstepsem}
	Let $\expr \in \exprs$ and $n \in \mathbb{N}$. We have
	\begin{enumerate}
		\item\label{lem:props:bigstepsem1} $\bigsemstep{\expr}{n}(\{\error\}) \leq \bigsemstep{\expr}{n+1}(\{\error\})$
		%
		\item\label{lem:props:bigstepsem2} $\bigsemstep{\expr}{n}(\{\obsfail\}) \leq \bigsemstep{\expr}{n+1}(\{\obsfail\})$
		%
		\item\label{lem:props:bigstepsem3} $\bigsemstep{\expr}{n}(\nonvals) \geq \bigsemstep{\expr}{n+1}(\nonvals)$.
	\end{enumerate}
\end{lemma}
%
\begin{proof}
	By induction on $n$, using that for $\expr' \in \{ \error,\obsfail,\val \}$, we have $\sem{\expr'} = \delta_{\expr'}$.
\end{proof}
%
%
Given $\expr\in\extexprs$ and $n\in\Nats$, define $\bigsemstepterm{\expr}{n} = \lambda \measset. \bigsemstep{\expr}{n}(\measset \setminus \nonvals)$. \kevin{submeasure}
\begin{definition}[Big-Step Semantics of Expressions]
	Let $\expr \in \exprs$. The \emph{big-step semantics of $\expr$} is 
	%
	\[
		\bigsem{\expr} ~{}={}~ \lim_{n \to \infty} \lambda \measset.
		%
		\bigsemstepterm{\expr}{n}~.
	\]\kevin{submeasure}
\end{definition}
%
The above limit is to be understood pointwise and exists by the monotone convergence theorem:  Since $\extexprs$ is the disjoint union 
%
\[
	 \exprs \discup  \{ \error \} \discup \{ \obsfail \} 
	 \quad {}={} \quad
	 \nonvals \discup \vals \discup \{ \error \} \discup \{ \obsfail \}~,
\]
%
every $\probmeas \in \distextexprs$ is the disjoint union of some measurable sets $\measset_{\nonvals} \subseteq \nonvals, \measset_{\vals}\subseteq\vals, \measset_{\error}\subseteq \{\error\},\measset_{\obsfail}\subseteq\{\obsfail\}$. 
\kevin{I'll have to think about this}
%
%\begin{align*}
%	\bigsemstepterm{\expr}{n}(\measset) = 
%	\bigsemstepterm{\expr}{n}(\nonvals) 
%	+\bigsemstepterm{\expr}{n}(\vals) 
%	+\bigsemstepterm{\expr}{n}(\measset) 
%	+\bigsemstepterm{\expr}{n}(\measset) 
%\end{align*}

\subsection{Measure Space on Expressions}

To make the probabilistic semantics precise, we need to define an appropriate measure space structure on the set of extended expressions. We construct this measure space using the notion of \emph{expression skeletons}.

\subsubsection{Expression Skeletons}

An expression skeleton is an extended expression with ``holes'' where real-valued constants can appear. Formally, we define:

\begin{definition}[Expression Skeleton]
An expression skeleton $s$ is an extended expression where each occurrence of a float constant is replaced by a hole $\square_i$ for some index $i \in \mathbb{N}$. We denote by $\text{Skel}$ the set of all expression skeletons, and by $\text{holes}(s)$ the number of distinct holes in skeleton $s$.
\end{definition}

For example, the expression $\uniform(0.0, 1.0) < 0.5$ has the skeleton $\uniform(\square_1, \square_2) < \square_3$ with three holes.

\subsubsection{Measure Space Construction}

Given a skeleton $s$ with $n = \text{holes}(s)$ holes, we can instantiate it with a vector $\vec{r} = (r_1, \ldots, r_n) \in \mathbb{R}^n$ to produce a concrete expression $s[\vec{r}]$ by replacing each hole $\square_i$ with the real value $r_i$.

The set of all extended expressions can then be represented as:
\[
\extexprs = \bigsqcup_{s \in \text{Skel}} \{s[\vec{r}] \mid \vec{r} \in \mathbb{R}^{\text{holes}(s)}\}
\]

This is a countable disjoint union because:
\begin{enumerate}
    \item The set of skeletons $\text{Skel}$ is countable (expressions have finite syntax trees)
    \item For each skeleton $s$, the set of instantiations forms a copy of $\mathbb{R}^{\text{holes}(s)}$
    \item Different skeletons produce disjoint sets of expressions
\end{enumerate}

\subsubsection{Measurable Sets}

We equip $\extexprs$ with the $\sigma$-algebra generated by sets of the form:
\[
\{s[\vec{r}] \mid \vec{r} \in B\}
\]
where $s$ is a skeleton and $B \subseteq \mathbb{R}^{\text{holes}(s)}$ is a Borel measurable set.

This construction ensures that:
\begin{itemize}
    \item Basic operations like evaluating an expression to a value are measurable
    \item The distributions returned by our probabilistic semantics are well-defined probability measures
    \item We can integrate over expressions when needed for semantic definitions
\end{itemize}

\subsubsection{Distributions Over Expressions}

When we write $\mathcal{D}(\extexprs)$ in our probabilistic semantics, we mean the space of probability measures on $(\extexprs, \mathcal{F})$ where $\mathcal{F}$ is the $\sigma$-algebra described above.

For continuous distributions like $\uniform(a, b)$, the implicit coercion to $\mathcal{D}(\extexprs)$ works as follows:
\begin{enumerate}
    \item The uniform distribution on $[a, b)$ is a measure on $\mathbb{R}$
    \item We lift it to expressions via the skeleton containing a single hole: $s = \square_1$
    \item The resulting distribution on expressions assigns probability to sets of the form $\{r \mid r \in B\}$ for Borel sets $B \subseteq [a, b)$
\end{enumerate}

This measure-theoretic foundation ensures that our probabilistic semantics is mathematically rigorous and that operations like conditioning (via observations) and marginalization are well-defined. 


\subsection{Defining $\mu_e$, $\mu_{e, \discretize{e}}$, $R_{\tau}$}
We consider a subset of our language, focusing only on let statements, conditionals, less than comparisons, and uniform and discrete distributions. 

% \begin{itemize}
% 	\item $\monunit \colon X \rightarrow \dist{X}$
%     \item $\monbind \colon \dist{X} \times (X \to \dist{Y}) \rightarrow \dist{Y}$
% \end{itemize}

\begin{definition}[Sub-Distribution]
	A sub-probability measure (the continuous analogue of a sub-distribution) on a measureable space $(X, \Sigma)$ is a measure $\mu: \Sigma \to [0,1]$ such that $\mu(X) \leq 1$. We write $D(X)$ for the set of all sub-probability measures of $X$. 
\end{definition}

\begin{definition}[Support]
	The support of $\mu \in \mathcal{D}(X)$ is the set of points $x \in X$ such that every open neighborhood of $x$ has positive measure:
	\[
	{supp}(\mu) \triangleq \{ x \in X \mid \forall U \ni x, \; U \text{ open}, \; \mu(U) > 0 \}.
	\]
\end{definition}

Recall the Maybe monad in Haskell:
\begin{align*}
	\texttt{return} &:: a \to \texttt{Maybe}\;a \\
	\texttt{return}\;x &= \texttt{Just}\;x \\[6pt]
	(\texttt{>>=}) &:: \texttt{Maybe}\;a \to (a \to \texttt{Maybe}\;b) \to \texttt{Maybe}\;b \\
	mx \;\texttt{>>=}\; f &= \texttt{case}\; mx\;\texttt{ of} \\
	&\quad \texttt{Nothing} \;\to\; \texttt{Nothing} \\
	&\quad \texttt{Just}\;x \;\to\; f\;x
\end{align*}

We consider a probabilistic analogue of this, defined below:

\begin{definition}[Probability Monad]

	\begin{enumerate}
	\item The return operator at $x \in X$ is the Dirac measure
	\begin{align*}
		\text{return} :: X \to \mathcal{D}(X + \bot) \\
		\text{return} \; x = \delta_{x} \text{ for every } x \in X 
	\end{align*}

	\item The bind operator is defined by
	\begin{align*}
		\text{bind} : \mathcal{D}(X + \bot) \;\to\; (X \to \mathcal{D}(Y + \bot)) \;\to\; \mathcal{D}(Y + \bot) \\
		\text{bind}(d, f) = \sum_{x \in support(d)} d(x) \cdot f_{\bot}(x) \\
		% \text{bind}(d,f) = \int_{x \in support(d)} f_{\bot}(x) \, d\,d(x)
	\end{align*}
	where $f_{\bot}(\bot) = \delta_{\bot}$, $f_{\bot}(x) = f(x) \text{ for every } x \in X$.
	\end{enumerate}
\end{definition}

\subsubsection{Small-step semantics for $e : \tau$. $Expr_{\tau} \to \mathcal{D}(Expr_{\tau} + \bot)$.}
We assume the expressions are well-typed. Since we assume all expressions are well-typed, we do not need to explicitly consider cases where type checking fails. Moreover, certain constructs in our language are only applicable to specific types.
{
\footnotesize
\begin{align*}
	\sem{ \letkw \; x = e_1 : \tau_1 \; \inkw \; e_2 : \tau_2 : \tau_2} &= \begin{cases}
		\delta_{e_2[v/x] : \tau_2} \quad &\text{ if } e_1 = v \\
		\sem{ e_1 : \tau_1 } \gg\!= \lambda f. \sem{ \letkw \; x = f \; \inkw \; e_2 : \tau_2 : \tau_2} \quad &\text{ otherwise } \\
	\end{cases} \\
	\sem{ \ifkw \; e_1 : \bool \; \thenkw \; e_2 : \tau \; \elsekw \; e_3 : \tau : \tau} &= \begin{cases}
		\delta_{\expr_2 : \tau}  \quad &\text{ if $e_1 = \val_1 = \true$} \\
		\delta_{\expr_3 : \tau}  \quad &\text{ if $e_1 = \val_1 = \false$} \\
		\sem{ e_1 : \bool} \gg\!= \lambda f. \sem{ \ifkw \; f \; \thenkw \; e_2 : \tau \; \elsekw \; e_3 : \tau : \tau} \quad &\text{ otherwise } \\
	\end{cases} \\
	\sem{ e_1 : \tau_1 < e_2 : \tau_2 : \bool} &= \begin{cases}
		\delta_{\true : \bool} \quad &\text{ if } e_1 = v_1, e_2 = v_2, v_1 < v_2 \\
		\delta_{\false : \bool} \quad &\text{ if } e_1 = v_1, e_2 = v_2, v_1 \geq v_2 \\
		\sem{ e_1 : \tau_1 } \gg\!= \lambda f. \sem{ f < v_2 : \tau_2 : \bool} \quad &\text{ if } e_2 = v_2 \\
		\sem{ e_2 : \tau_2 } \gg\!= \lambda f. \sem{ v_1 : \tau_1 < f : \bool} \quad &\text{ if } e_1 = v_1 \\
		\sem{ e_1 : \tau_1 } \gg\!= \lambda f. \sem{ f < e_2 : \tau_2 : \bool} \quad &\text{ otherwise } \\
	\end{cases} \\
	&\text{where } \tau_1 = \tau_2 = \floattype{B}{V} \\
	\sem{ e_1 : \tau_1 \finlt{n} e_2 : \tau_2 : \bool} &= \begin{cases}
		\delta_{\true : \bool} \quad &\text{ if } e_1 = k_{1,\#n}, e_2 = k_{2,\#n}, k_1 < k_2 \\
		\delta_{\false : \bool} \quad &\text{ if } e_1 = k_{1,\#n}, e_2 = k_{2,\#n}, k_1 \geq k_2 \\
		\sem{ e_1 : \tau_1 } \gg\!= \lambda f. \sem{ f \finlt{n} k_{2,\#n} : \tau_2 : \bool} \quad &\text{ if } e_2 = k_{2,\#n} \\
		\sem{ e_2 : \tau_2 } \gg\!= \lambda f. \sem{ k_{1,\#n} : \tau_1 \finlt{n} f : \bool} \quad &\text{ if } e_1 = k_{1,\#n} \\
		\sem{ e_1 : \tau_1 } \gg\!= \lambda f. \sem{ f \finlt{n} e_2 : \tau_2 : \bool} \quad &\text{ otherwise } \\
	\end{cases} \\
	&\text{where } \tau_1 = \tau_2 = \fin{n} \\
	\sem{ \uniform(e_1 : \tau_1, e_2 : \tau_2) : \tau } &= \begin{cases}
		\text{Uniform}{(v_1, v_2)} \quad &\text{ if $e_1 = v_1$, $e_2 = v_2$, $v_1 \leq v_2$} \\
		\delta_{\bot} \quad &\text{ if $e_1 = v_1$, $e_2 = v_2$, $v_1 > v_2$} \\
		\sem{ e_1 : \tau_1 } \gg\!= \lambda f. \sem{ \uniform(f, v_2 : \tau_2) : \tau } \quad &\text{ if } e_2 = v_2 \\
		\sem{ e_2 : \tau_2 } \gg\!= \lambda f. \sem{ \uniform(v_1 : \tau_1, f) : \tau } \quad &\text{ if } e_1 = v_1 \\
		\sem{ e_1 : \tau_1 } \gg\!= \lambda f. \sem{ \uniform(f, e_2 : \tau_2) : \tau } \quad &\text{ otherwise } \\
	\end{cases} \\
	&\text{where } \tau = \floattype{B}{V} \\
	\sem{ \discrete(v_0, \ldots, v_n) : \tau } &= \begin{cases}
		\text{Discrete}{(v_0, \ldots, v_n)} \quad &\text{ if } \sum_i v_i = 1 \\
		\delta_{\bot} \quad &\text{ otherwise }
	\end{cases} \\
	&\text{where } \tau = \fin{n} \\
	\sem{ v : \tau } &= \delta_{v:\tau} \\
	&\text{where } \tau = \floattype{B}{V} \\
	\sem{ \finconst{k}{n} : \tau } &= \delta_{\finconst{k}{n} : \tau} \\
	&\text{where } \tau = \fin{n} \\
\end{align*}
}

\subsubsection{Relation $\relation_{\tau}$}
We couple the continuous (and type-annotated) program $e$ appropriately with its discretized one $\discretize{e}$. Formally, the coupling relation is:

\[
\relation_{\tau}(e : \tau,e' : \tau') \iff e' : \tau' == \discretize{e : \tau} : \discretize{\tau'}.
\]

% What does this mean precisely? In particular, when the return type of $e$ is $\bool$, then the return type of $\discretize{e}$ is also $\bool$, and they have the same distribution. When the return type of $e$ is $\float$, then the return type of $\discretize{e}$ is also $\float$, and they have the same distribution. 

% \begin{align*}
% 	\{(\true,\true), (\false,\false)\} \in \relation_{\bool} \\
% 	\{(v,v) : v \in supp(\mu)\} \in \relation_{\float} \\
% 	\text{ where $\mu$ is the coupled distribution}
% \end{align*}

\subsubsection{Small-step semantics for $e : \tau, e' : \tau'$.  $\relation_{\tau} \to \mathcal{D}(\relation_{\tau} + \bot)$.} 
We assume the expressions $e, e'$ are well-typed. 

% \begin{definition}[Probability Monad]
% 	\begin{enumerate}
% 		\item The return operator is defined by
% 		\begin{align*}
% 			\text{return} :: (X \times X) \to \mathcal{D}((X + \bot) \times (X + \bot)) \\
% 			\text{return} \; (x,y) = \delta_{(x,y)} \text{ for every } (x,y) \in X \times X
% 		\end{align*}
	
% 		\item The bind operator is defined by
% 		\begin{align*}
% 			\text{bind} : \mathcal{D}((X + \bot) \times (X + \bot)) \;\to\; ((X \times X) \to \mathcal{D}(Y + \bot) \times (Y + \bot)) \;\to\; \mathcal{D}((Y + \bot) \times (Y + \bot)) \\
% 			\text{bind}(d, f) = \sum_{(x,y) \in support(d)} d(x,y) \cdot f_{\bot}(x,y) \\
% 		\end{align*}
		
% 		where: 
% 		$f_{\bot}(x,y) =
% 		\begin{cases}
% 			f(x,y) &\text{ if } x \in X, y \in X \\
% 			\delta_{(\pi_1(f(x,y)),\bot)} &\text{ if } x \in X, y = \bot \\
% 			\delta_{(\bot,\pi_2(f(x,y)))} &\text{ if } x = \bot, y \in X \\
% 			\delta_{(\bot,\bot)} &\text{ if } x = \bot, y = \bot \\
% 		\end{cases}$
% 		\end{enumerate}
% \end{definition}

{\tiny
\allowdisplaybreaks
\begin{align*}
	%% let
	\sem{ \letkw \; x = e_1 : \tau_1 \; \inkw \; e_2 : \tau_2 : \tau_2, \\[-3ex]
	\letkw \; x' = e_1' : \tau_1' \; \inkw \; e_2' : \tau_2' : \tau_2'} &= \begin{cases}
		\delta_{(e_2[v/x] : \tau_2, e_2'[v'/x'] : \tau_2')} \quad &\text{ if } e_1 = v, e_1' = v' \\
		\sem{ e_1 : \tau_1, e_1' : \tau_1'} \gg\!= \lambda (f, f'). \sem{ \letkw \; x = f \; \inkw \; e_2 : \tau_2 : \tau_2, \letkw \; x' = f' \; \inkw \; e_2' : \tau_2' : \tau_2'} \quad &\text{ if } e_1 \neq v, e_1' \neq v' \\
	\end{cases} \\
	%% if-else
	\sem{ \ifkw \; e_1 : \bool \; \thenkw \\
	\; e_2 : \tau \; \elsekw \; e_3 : \tau : \tau, \\
	\ifkw \; e_1' : \bool \; \thenkw \\[-5ex]
	\; e_2' : \tau' \; \elsekw \; e_3' : \tau' : \tau' } &= \begin{cases}
		\delta_{(e_2 : \tau, e_2' : \tau')}  \quad &\text{ if $e_1 = v = \true$, $e_1' = v' = \true$} \\
		\delta_{(e_3 : \tau, e_3' : \tau')}  \quad &\text{ if $e_1 = v = \false$, $e_1' = v' = \false$} \\
		\sem{ e_1 : \bool, e_1' : \bool } \gg\!= \lambda (f,f'). \sem{ \ifkw \; f \; \thenkw \; e_2 : \tau \; \elsekw \; e_3 : \tau : \tau, \\
		\ifkw \; f' \; \thenkw \; e_2' : \tau' \; \elsekw \; e_3' : \tau' : \tau'} \quad &\text{ if } e_1 \neq v, e_1' \neq v' \\
	\end{cases} \\
	%% less than - float
	\sem{ e_1 : \tau_1 < e_2 : \tau_2 : \bool, \\[-5ex]
	e_1' : \tau_1 < e_2' : \tau_2 : \bool } &= \begin{cases}
		\delta_{(v_1 < v_2 : \bool, v_1' < v_2' : \bool)} \quad &\text{ if } e_1 = v_1, e_1' = v_1', e_2 = v_2, e_2' = v_2' \\
		\sem{ e_1 : \tau_1, e_1' : \tau_1} \gg\!= \lambda (f,f'). \sem{ f < v_2 : \tau_2, f' < v_2 : \tau_2 } \quad &\text{ if } e_1 \neq v_1, e_1' \neq v_1', e_2 = e_2' = v_2 \\
		\sem{ e_2 : \tau_2, e_2' : \tau_2 } \gg\!= \lambda (f,f'). \sem{ v_1 : \tau_1 < f, v_1 : \tau_1 < f' } \quad &\text{ if } e_1 = e_1' = v_1, e_2 \neq v_2, e_2' \neq v_2' \\
		\sem{ e_1 : \tau_1, e_1' : \tau_1 } \gg\!= \lambda (f,f'). \sem{ f < e_2 : \tau_2, f' < e_2' : \tau_2 } \quad &\text{ if } e_1 \neq v_1, e_1' \neq v_1', e_2 \neq v_2, e_2' \neq v_2' \\
	\end{cases} \\
	&\text{where } \tau_1 = \tau_2 = \floattype{B}{V}, B = \top \\
	%% less than - fin
	\sem{ e_1 : \tau_1 < e_2 : \tau_2 : \bool, \\[-5ex]
	e_1' : \tau_1' \finlt{n+1} e_2' : \tau_2' : \bool } &= \begin{cases}
		\delta_{(v_1 < v_2 : \bool, k_1 < k_2 : \bool)} \quad &\text{ if } e_1 = v_1, e_1' = k_{1,\#n+1}, e_2 = v_2, e_2' = k_{2,\#n+1} \\
		\sem{ e_1 : \tau_1, e_1' : \tau_1' } \gg\!= \lambda (f,f'). \sem{ f < v_2 : \tau_2, f' \finlt{n+1} k_{2,\#n+1} : \tau_2' } \quad &\text{ if } e_1 \neq v_1, e_1' \neq k_{1,\#n+1}, e_2 = v_2, e_2' = k_{2,\#n+1} \\
		\sem{ e_2 : \tau_2, e_2' : \tau_2' } \gg\!= \lambda (f,f'). \sem{ v_1 : \tau_1 < f, k_{1,\#n+1} : \tau_1' \finlt{n+1} f' } \quad &\text{ if } e_1 = v_1, e_1' = k_{1,\#n+1}, e_2 \neq v_2, e_2' \neq k_{2,\#n+1} \\
		\sem{ e_1 : \tau_1, e_1' : \tau_1' } \gg\!= \lambda (f,f'). \sem{ f < e_2 : \tau_2, f' \finlt{n+1} e_2' : \tau_2' } \quad &\text{ if } e_1 \neq v_1, e_1' \neq k_{1,\#n+1}, e_2 \neq v_2, e_2' \neq k_{2,\#n+1} \\
	\end{cases} \\
	&\text{where } \tau_1 = \tau_2 = \floattype{B}{V}, B = \{\sim_1\!\!b_1, \ldots, \sim_n\!\!b_n\}, \tau_1' = \tau_2' = \fin{n+1} \\
	%% uniform uniform
	\sem{ \uniform(e_1 : \tau_1, e_2: \tau_2) : \tau, \\[-5ex]
	\uniform(e_1' : \tau_1, e_2' : \tau_2) : \tau} &= \begin{cases}
		\text{Uniform}(v_1, v_2) \gg\!= \lambda v.\delta_{(v, v)} \quad &\text{ if } e_1 = e_1' = v_1 \leq e_2 = e_2' = v_2 \\ 
		\delta_{\bot} \quad &\text{ if } e_1 = v_1 > e_2 = v_2 \\
		\sem{ e_1 : \tau_1, e_1' : \tau_1' } \gg\!= \lambda (f,f'). \sem{ \uniform(f, v_2 : \tau_2) : \tau, \uniform(f',v_2 : \tau_2) : \tau } \quad &\text{ if } e_1 \neq v_1, e_1' \neq v_1', e_2 = e_2' = v_2 \\
		\sem{ e_2 : \tau_2, e_2' : \tau_2' } \gg\!= \lambda (f,f'). \sem{ \uniform(v_1 : \tau_1, f) : \tau, \uniform(v_1 : \tau_1, f') : \tau } \quad &\text{ if } e_1 = e_1' = v_1, e_2 \neq v_2, e_2' \neq v_2',\\
		\sem{ e_1 : \tau_1, e_1' : \tau_1' } \gg\!= \lambda (f,f'). \sem{ \uniform(f, e_2 : \tau_2) : \tau, \uniform(f',e_2' : \tau_2) : \tau } \quad &\text{ if } e_1 \neq v_1, e_1' \neq v_1', e_2 \neq v_2, e_2' \neq v_2' \\
	\end{cases} \\
	&\text{where } \tau = \floattype{B}{V}, B = \top \\
	%% uniform discrete
	\sem{ \uniform(e_1 : \tau_1, e_2 : \tau_2) : \tau, \\[-5ex]
	\discrete(v_0', \ldots, v_n') : \tau'} &= \begin{cases}
		\text{Uniform}(v_1, v_2) \gg\!= \lambda v.\delta_{(v, i)} \text{ s.t. } v \in \intervals{B}_i \quad &\text{ if } e_1 = v_1 \leq e_2 = v_2, \sum_i v_i' = 1, \\
		&v_i' = \text{CDF}_{\text{unif}(v_1,v_2)}(b_{i+1}) - \text{CDF}_{\text{unif}(v_1,v_2)}(b_{i}) \\
		\delta_{\bot} \quad &\text{ if } e_1 = v_1 > e_2 = v_2 \\
	\end{cases} \\
		&\text{ where } \tau = \floattype{B}{V}, B = \{\sim_1\!\!b_1, \ldots, \sim_n\!\!b_n\} \neq \top, \tau' = \fin{n+1} \\
	%% uniform discretecase
	\sem{ \uniform(e_1 : \tau_1, e_2 : \tau_2) : \tau, \\[-5ex]
	\texttt{DiscreteCase}(e_1' : \tau_1', e_2' : \tau_2') : \tau'} &= \begin{cases}
		\sem{e_1, e_1' } \gg\!= \lambda(u_1, i).\sem{ e_2, e_2' } \gg\!= \lambda(u_2, j).\begin{cases}
			% \sem{ \uniform(v_1, w_1) : \tau, \discrete(v_{000}, v_{001}, \ldots, v_{00n})} \\
			% \text{ if } u_1 = v_1, u_1' = 0, u_2 = w_1, u_2' = 0 \\
			% \ldots \\
			\sem{ \uniform(u_1, u_2) : \tau, \\
			\discrete(v_{ij0}, \ldots, v_{ijn}) : \tau'} \\
			% \ldots \\
			% \sem{ \uniform(v_r, w_s) : \tau, \discrete(v_{rs0}, v_{rs1}, \ldots, v_{rsn})} \\
			% \text{ if } u_1 = v_r, u_1' = l, u_2 = w_s, u_2' = m \\
		\end{cases} \\
		\text{s.t. } v_{ijk} = \text{CDF}_{\text{unif}(u_1, u_2)}(b_{k+1}) - \text{CDF}_{\text{unif}(u_1, u_2)}(b_{k}) \quad &\text{if } e_1 \neq v_1 \text{ or } e_2 \neq v_2 \\
	\end{cases} \\
	\texttt{DiscreteCase}(e_1,e_2) &= \letkw \; v_1' = e_1 \; \inkw \\
			&\letkw \; v_2' = e_2 \; \inkw \\
			&\ifkw\; v_1' \fineq{l+1} \finconst{0}{l+1}\; \logand \; v_2' \fineq{m+1} \finconst{0}{m+1}\; \thenkw \\
			&\quad \discrete(v_{000}, v_{001}, \ldots, v_{00n}) \\
			&\vdots \\
			&\elsekw\; \ifkw\; v_1' \fineq{l+1} \finconst{i}{l+1}\; \logand \; v_2' \fineq{m+1} \finconst{j}{m+1}\; \thenkw \\
			&\quad \discrete(v_{ij0}, v_{ij1}, \ldots, v_{ijn}) \\
			&\vdots \\
			&\elsekw \\
			&\quad \discrete(v_{rs0}, v_{rs1}, \ldots, v_{rsn}) \\
	&\text{where } \tau = \floattype{B}{V}, B = \{\sim_1\!\!b_1, \ldots, \sim_n\!\!b_n\} \neq \top, \\
	&\tau_1 = \floattype{B_1}{V_1}, B_1 = \{\sim_1\!\!c_1, \ldots, \sim_n\!\!c_l\}, V_1 = \{v_1, \ldots, v_r\}, \\
	&\tau_2 = \floattype{B_2}{V_2}, B_2 = \{\sim_1\!\!d_1, \ldots, \sim_n\!\!d_m\}, V_2 = \{w_1, \ldots, w_s\} \\
	&\tau' = \fin{n+1}, \tau_1' = \fin{l+1}, \tau_2' = \fin{m+1} \\
	%% value float
	\sem{ v : \tau, v' : \tau } &= \delta_{(v, v')} \text{ if } v = v' \\
	&\text{where } \tau = \floattype{B}{V}, B = \top \\
	%% value fin
	\sem{ v : \tau, \finconst{k}{n} : \tau' } &= \delta_{(v, \finconst{k}{n})} \text{ if } v \in \intervals{B}_k \\
	&\text{where } \tau = \floattype{B}{V}, B = \{\sim_1\!\!b_1, \ldots, \sim_n\!\!b_n\}, \tau' = \fin{n+1} \\
\end{align*}
}

\subsubsection{Lemmas}
We want to show the following lemmas for the pairing semantics:
\begin{lemma}[Progress with Error]
	If $(e_1 : \tau, e_1' : \tau')$, then either:
	\begin{enumerate}
		\item $(e_1, e_1')$ is a value pair
		\item There exists $(e_2, e_2')$ such that $\sem{e_1,e_1'}(\{(e_2,e_2')\}) = 1$
		\item $(e_1,e_1')$ is $\bot$
	\end{enumerate}
\end{lemma}

\begin{lemma}[Preservation]
	If $(e_1 : \tau, e_1' : \tau')$ and $\sem{e_1,e_1'}(\{(e_2,e_2')\}) = 1$ where $(e_2,e_2') \neq \top$, then $(e_2 : \tau, e_2' : \tau')$.
\end{lemma}

\begin{lemma}[Substitution]
	If $\Gamma, (x : \tau_1, x' : \tau_1') \vdash (e_2 : \tau_2, e_2' : \tau_2')$ and $\Gamma \vdash (v : \tau_1, v' : \tau_1')$, then $\Gamma \vdash (e_2[v/x] : \tau_2, e_2'[v'/x'] : \tau_2')$.
\end{lemma}

\begin{lemma}
	If $\relation(e,e')$ holds, then $\sem{ e,e' }(\{(f,f') : R(f,f')\}) = 1$. 
\end{lemma}

From the above lemmas, we can conclude that $e$ and $\discretize{e}$ are semantically equivalent, i.e., they produce the same output distribution up to discretization of float values. 


% of two expressions executing in lockstep. We run both programs side-by-side under the coupled small-step rules and relate $(e, \discretize{e})$ by $R_{\tau}$ at every step of the execution. At termination, both programs produce values related by $R_{\tau}$, which means their distributions agree.

% Configurations are of the form $(e, t, p)$, where:
% \begin{itemize}
% 	\item $e$ is the expression
% 	\item $t$ is the trace (sequence of samples)
% 	\item $p$ is the current probability weight
% 	\item \textbf{Initial configuration:} $(e, \varepsilon, 1.0)$ where $\varepsilon$ is the empty trace
% \end{itemize}

% \subsubsection{Small-step semantics for $\mu_e$}
% \begin{mathpar}
% 	\inferrule[\textsc{Let}]
% 	{\ }
% 	{(\letkw \; x = v \; \inkw \; e, t, p) \rightarrow (\{e[v/x]\}, t, p)}

% 	\inferrule[\textsc{IfTrue}]
% 	{\ }
% 	{(\ifkw \; \text{true} \; \thenkw \; e_1 \; \elsekw \; e_2, t, p) \rightarrow (\{e_1\}, t, p)}
	
% 	\inferrule[\textsc{IfFalse}]
% 	{\ }
% 	{(\ifkw \; \text{false} \; \thenkw \; e_1 \; \elsekw \; e_2, t, p) \rightarrow (\{e_2\}, t, p)}

% 	\inferrule[\textsc{LessTrue}]
% 	{c_1 < c_2}
% 	{(c_1 < c_2, t, p) \rightarrow (\{\text{true}\}, t, p)}

% 	\inferrule[\textsc{LessFalse}]
% 	{c_1 \geq c_2}
% 	{(c_1 < c_2, t, p) \rightarrow (\{\text{false}\}, t, p)}

% 	\inferrule[\textsc{Uniform}]
% 	{v_1, v_2 \text{ are float values} \\ v_1 < v_2}
% 	{(\uniform(v_1, v_2), t, p) \rightarrow (c, t::c, p * (1/(v_2-v_1)))}

% 	\inferrule[\textsc{Discrete}]
% 	{p_0, \ldots, p_n \text{ are probabilities} \\ \sum_{i=0}^{n} p_i = 1}
% 	{(\discrete(p_0, \ldots, p_n), t, p) \rightarrow (i, t::i, p * p_i)}
% \end{mathpar}

% \subsubsection{Small-step semantics for $\mu_{e, \discretize{e}}$}
% \begin{mathpar}
% 	\inferrule[\textsc{Let-2}]
% 	{e := \letkw \; x_1 = v_1 \; \inkw \; e_1 \ e' := \letkw \; x_2 = v_2 \; \inkw \; e_2}
% 	{((e, t, p), (e', t', p')) \rightarrow ((e_1[v_1/x_1], t, p), (e_2[v_2/x_2], t', p'))}
	
% 	\inferrule[\textsc{IfTrue-2}]
% 	{e := \ifkw \; \text{true} \; \thenkw \; e_1 \; \elsekw \; e_2 \\ e' := \ifkw \; \text{true} \; \thenkw \; e_1' \; \elsekw \; e_2'}
% 	{((e, t, p), (e', t', p')) \rightarrow ((e_1, t, p), (e_1', t', p'))}
	
% 	\inferrule[\textsc{IfFalse-2}]
% 	{e := \ifkw \; \text{false} \; \thenkw \; e_1 \; \elsekw \; e_2 \\ e' := \ifkw \; \text{false} \; \thenkw \; e_1' \; \elsekw \; e_2'}
% 	{((e, t, p), (e', t', p')) \rightarrow ((e_2, t, p), (e_2', t', p'))}
	
% 	\inferrule[\textsc{LessTrue-2}]
% 	{c_1 < c_2 \\ c_1' < c_2'}
% 	{((c_1 < c_2, t, p), (c_1' < c_2', t', p')) \rightarrow ((\text{true}, t, p), (\text{true}, t', p'))}
	
% 	\inferrule[\textsc{LessFalse-2}]
% 	{c_1 \geq c_2 \\ c_1' \geq c_2'}
% 	{((c_1 < c_2, t, p), (c_1' < c_2', t', p')) \rightarrow ((\text{false}, t, p), (\text{false}, t', p'))}
	
% 	\inferrule[\textsc{Uniform-Discrete-2}]
% 	{v_1, v_2 \text{ are float values} \\ v_1 < v_2 \\ p_0, \ldots, p_n \text{ are probabilities} \\ c \sim \text{Uniform}(v_1, v_2) \\ i \sim \text{Discrete}(p_0, \ldots, p_n)}
% 	{((\uniform(v_1, v_2), t, p), (\discrete(p_0, \ldots, p_n), t', p')) \rightarrow ((c, t::c, p \times \frac{1}{v_2 - v_1}), (i, t'::i, p' \times p_i))}
	
% 	\inferrule[\textsc{Uniform-Uniform-2}]
% 	{v_1, v_2 \text{ are float values} \\ v_1 < v_2 \\ v_1', v_2' \text{ are float values} \\ v_1' < v_2' \\ c \sim \text{Uniform}(v_1, v_2) \\ c' \sim \text{Uniform}(v_1', v_2')}
% 	{((\uniform(v_1, v_2), t, p), (\uniform(v_1', v_2'), t', p')) \rightarrow ((c, t::c, p \times \frac{1}{v_2 - v_1}), (c', t'::c', p' \times \frac{1}{v_2' - v_1'}))}

% \end{mathpar}

% \subsubsection{Relation $R_{\tau}(e, \discretize{e})$}

% \begin{align*}
% 	\mathcal{R_{\bool}}(e, \discretize{e}) &\iff \{(\true, \true), (\false, \false)\} \\
% 	\mathcal{R}_{\floattype{B}{V}}(e, \discretize{e}) &\iff 
% 	\begin{cases}
% 		\discretize{e} = 0 \text{ if } e \in \intervals{B}_0 \\
% 		\discretize{e} = 1 \text{ if } e \in \intervals{B}_1 \\
% 		\ldots \\
% 		\discretize{e} = n \text{ if } e \in \intervals{B}_n \\
% 	\end{cases} \text{where } B = \{b_1, \ldots, b_n\} \\
% 	\mathcal{R}_{\floattype{B}{V}}(e, \discretize{e}) &\iff e = \discretize{e} \text{ where } B = \top \\
% 	\mathcal{R}_{\fin{n}}(e, \discretize{e}) &\iff e = \discretize{e} \\
% 	\mathcal{R}_{\intty}(e, \discretize{e}) &\iff e = \discretize{e} \\
% 	\mathcal{R}_{\tau_1 \times \tau_2}(e, \discretize{e}) &\iff (\pi_1(e), \pi_1(\discretize{e})) \in \mathcal{R}_{\tau_1} \land
% 	(\pi_2(e), \pi_2(\discretize{e})) \in \mathcal{R}_{\tau_2} \\
% 	\mathcal{R}_{\tau_1 \to \tau_2}(e, \discretize{e})
%     &\iff \forall (v, v') \in \mathcal{R}_{\tau_1}, ((\lambda x : \tau_1. e) \; v, (\lambda x : \tau_1. \discretize{e}) \; v') \in \mathcal{R}_{\tau_2} \\
% 	\mathcal{R}_{\listty{\tau}}(e, \discretize{e}) &\iff 
% 	\begin{cases}
% 		e = \text{nil} \land \discretize{e} = \text{nil} \\
% 		e = v::vs \land \discretize{e} = v'::vs' \land (v, v') \in \mathcal{R}_{\tau} \land (vs, vs') \in \mathcal{R}_{\listty{\tau}}
% 	\end{cases}
% \end{align*}

% \subsection{Lemmas}

% \begin{enumerate}
% 	\item $\mu_{e, \discretize{e}}(R_{\tau}) = 1$
% 	\item $\pi_1(\mu_{e, \discretize{e}}) = \mu_e$, $\pi_2(\mu_{e, \discretize{e}}) = \mu_{\discretize{e}}$
% \end{enumerate}

% \subsection{Example}
% \begin{figure}[h]
% 	\centering
% 	\begin{subfigure}[t]{0.5\textwidth}
% 	\scriptsize % shrink so it fits
% 	\begin{align*}
% 		&(\texttt{let x = uniform(0,2) in} \\
% 		&\;\;\;\texttt{if x < 0.5 then uniform(0,2)} \\
% 		&\;\;\;\texttt{else uniform(0,4)}, \varepsilon, 1.0) \\
% 		&\rightarrow (\texttt{if uniform(0,2) < 0.5 then} \\
% 		&\;\;\;\texttt{uniform(0,2) else uniform(0,4)}, \varepsilon, 1.0) \quad \textsc{Let} \\
% 		&\rightarrow (\texttt{if 0.3 < 0.5 then uniform(0,2)} \\
% 		&\;\;\;\texttt{else uniform(0,4)}, 0.3, 0.5) \quad \textsc{Uniform} \\
% 		&\rightarrow (\texttt{if \true then uniform(0,2)} \\
% 		&\;\;\;\texttt{else uniform(0,4)}, 0.3, 0.5) \quad \textsc{LessTrue} \\
% 		&\rightarrow (\texttt{uniform(0,2)}, 0.3, 0.5) \quad \textsc{IfTrue} \\
% 		&\rightarrow (0.2, 0.3::0.2, 1/4) \quad \textsc{Uniform}
% 	\end{align*}
% 	\caption{Continuous version}
% 	\end{subfigure}\hfill
% 	\begin{subfigure}[t]{0.5\textwidth}
% 	\scriptsize
% 	\begin{align*}
% 		&(\texttt{let x = discrete(0.25, 0.75) in} \\
% 		&\;\;\;\texttt{if x < 1 then uniform(0,2)} \\
% 		&\;\;\;\texttt{else uniform(0,4)}, \varepsilon, 1.0) \\
% 		&\rightarrow (\texttt{if discrete(0.25, 0.75) < 1 then} \\
% 		&\;\;\;\texttt{uniform(0,2) else uniform(0,4)}, \varepsilon, 1.0) \quad \textsc{Let} \\
% 		&\rightarrow (\texttt{if 0 < 1 then uniform(0,2)} \\
% 		&\;\;\;\texttt{else uniform(0,4)}, 0, 0.25) \quad \textsc{Discrete} \\
% 		&\rightarrow (\texttt{if \true then uniform(0,2)} \\
% 		&\;\;\;\texttt{else uniform(0,4)}, 0, 0.25) \quad \textsc{LessTrue} \\
% 		&\rightarrow (\texttt{uniform(0,2)}, 0, 0.25) \quad \textsc{IfTrue} \\
% 		&\rightarrow (0.2, 0::0.2, 1/8) \quad \textsc{Uniform}
% 	\end{align*}
% 	\caption{Discrete version}
% 	\end{subfigure}
% \end{figure}


% \begin{align*}
% 	&(\texttt{let x = uniform(0,2) in if x < 0.5 then uniform(0,2) else uniform(0,4)}, \varepsilon, 1.0) \\
% 	&\rightarrow (\texttt{if uniform(0,2) < 0.5 then uniform(0,2) else uniform(0,4)}, \varepsilon, 1.0) \quad \textsc{Let} \\
% 	&\rightarrow (\texttt{if 0.3 < 0.5 then uniform(0,2) else uniform(0,4)}, 0.3, 0.5) \quad \textsc{Uniform} \\
% 	&\rightarrow (\texttt{if \true then uniform(0,2) else uniform(0,4)}, 0.3, 0.5) \quad \textsc{LessTrue} \\
% 	&\rightarrow (\texttt{uniform(0,2)}, 0.3, 0.5) \quad \textsc{IfTrue} \\
% 	&\rightarrow (0.2, 0.3::0.2, 1/4) \quad \textsc{Uniform} \\
% \end{align*}

% \begin{align*}
% 	&(\texttt{let x = discrete(0.25, 0.75) in if x < 1 then uniform(0,2) else uniform(0,4)}, \varepsilon, 1.0) \\
% 	&\rightarrow (\texttt{if discrete(0.25, 0.75) < 1 then uniform(0,2) else uniform(0,4)}, \varepsilon, 1.0) \quad \textsc{Let} \\
% 	&\rightarrow (\texttt{if 0 < 1 then uniform(0,2) else uniform(0,4)}, 0, 0.25) \quad \textsc{Discrete} \\
% 	&\rightarrow (\texttt{if \true then uniform(0,2) else uniform(0,4)}, 0, 0.25) \quad \textsc{LessTrue} \\
% 	&\rightarrow (\texttt{uniform(0,2)}, 0, 0.25) \quad \textsc{IfTrue} \\
% 	&\rightarrow (0.2, 0::0.2, 1/8) \quad \textsc{Uniform} \\
% \end{align*}





% \subsection{Contextual Equivalence}
% We give a defintion of contextual equivalence below. Two expressions are contextually equivalent ($=_{ctx}$) if for all closing program contexts $C$, their observable aggregate behavior is the same. We take programs to be real values closed expressions, and their observable behavior consists of their measures ($\Sigma_{\R} \rightarrow \R^{+}$).


% We give the correctness statement for our discretization; we want the original expression $e$ and the discretized expression $\discretize{e}$ to be contextually equivalent. Two expressions are contextually equivalent ($=_{ctx}$) if for all closing program contexts $C$ their observable aggregate behavior is the same. We take programs to be real values closed expressions, and their observable behavior consists of their measures ($\Sigma_{\R} \rightarrow \R^{+}$).

% \begin{definition}[Contextual Equivalence / Correctness]
% 	If $\Gamma \vdash e_1 : \tau_1$ and $\Gamma \vdash e_2 : \tau_2$, then $e_1$ and $e_2$ are contextually equivalent $(e_1 =_{ctx} e_2)$ if and only if for all contexts $C$ such that $C : \Gamma$ and $\vdash C[e_1] : \float$ and $\vdash C[e_2] : \float$ (or $\vdash C[e_1] : \bool$ and $\vdash C[e_2] : \bool$) and for all measurable sets $A \in \Sigma_{\R}$,

% 	\[
% 	\mu_{C[e_1]}(A) = \mu_{C[e_2]}(A).
% 	\]
% \end{definition}

% $C : \Gamma$ means that $C$ provides bindings satisfying $\Gamma$ to the expression placed in its hole. 

% \subsection{Logical Relation for Contextual Equivalence}
% We develop a logical relation for proving expressions contextually equivalent. Figure~\ref{fig:logical-relation} gives the logical relation and its auxillary relations.

% \begin{figure}
% 	\begin{align*}
% 		\Gamma \vdash e_1 \approx e_2 : \tau 
%         &\iff \forall (\gamma_1, \gamma_2) \in \mathcal{G}\sem{\Gamma}, (e_1 \cdot \gamma_1, e_2 \cdot \gamma_2) \in \mathcal{E}\sem{\tau} \\
% 		(r_1, r_2) \in \mathcal{V}\sem{\tau} 
%         &\iff r_1 = r_2 \\[1ex]
%         (\lambda x.\, e_1, \lambda x.\, e_2) \in \mathcal{V}\sem{\tau_1 \to \tau_2} 
%         &\iff \forall (v_1, v_2) \in \mathcal{V}\sem{\tau_1},\; (e_1[x \mapsto v_1], e_2[x \mapsto v_2]) \in \mathcal{E}\sem{\tau_2} \\
%         (\gamma_1, \gamma_2) \in \mathcal{G}\sem{\Gamma} 
%         &\iff \gamma_1 \models \Gamma \;\wedge\; \gamma_2 \models \Gamma \;\wedge\; \forall x \in \mathrm{dom}(\Gamma), (\gamma_1(x), \gamma_2(x)) \in \mathcal{V}\sem{\Gamma(x)} \\
%         (e_1, e_2) \in \mathcal{E}\sem{\tau} 
%         &\iff \forall (A_1, A_2) \in \mathcal{A}\sem{\tau},\; \mu_{e_1}(A_1) = \mu_{e_2}(A_2) \\
%         (A_1, A_2) \in \mathcal{A}\sem{\tau} 
%         &\iff A_1, A_2 \in \Sigma\sem{\tau} \;\wedge\; \forall (v_1, v_2) \in \mathcal{V}\sem{\tau},\; (v_1 \in A_1 \iff v_2 \in A_2)
% 	\end{align*}
% 	\caption{Logical relation and auxiliary relations. It defines the relation $\Gamma \vdash e_1 \approx e_2 : \tau$.}
% 	\label{fig:logical-relation}
% \end{figure}

% \begin{itemize}
% 	\item $\mathcal{V}\sem{\tau}$ relates closed values of type $\tau$. Real values are related if they are identitcal, and functions are related if they take related inputs to related evaluation configurations.
% 	\item $\mathcal{G}\sem{\Gamma}$ relates substitutions.
% 	\item $\mathcal{E}\sem{\tau}$ relates closed expressions. Expressions are related if their measures agree on emasurable sets related by $\mathcal{A}\sem{\tau}$.
% 	\item $\mathcal{A}\sem{\tau}$ relates measurable sets of values.
% \end{itemize}

\section{Code Sketches for Synthetic Benchmarks}\label{appen:synthetic-sketches}
This section contains code snippets of the programs used in the scaling experiments in Section~\ref{sec:synthetic-benchmarks}. We implemented a Python script that automatically generates these programs based on configurable parameters such as program size, guard structure, and dependency pattern. For reproducibility, we used a fixed random seed of random.seed(11). All $r_i$ values in the generated programs, where $0 \leq r_i \leq 1$ is randomly chosen, are rounded to two decimal places.


\subsection{Conditional Independent Programs}
\begin{figure}
    \begin{lstlisting}
        let x1 = uniform(0,1) in
        let x2 = if x1 < 0.56 then uniform(0,2) else uniform(0,3) in
        let x3 = if x2 < 0.92 then uniform(0,4) else uniform(0,5) in
        let x4 = if x3 < 0.47 then uniform(0,6) else uniform(0,7) in
        let x5 = if x4 < 0.51 then uniform(0,8) else uniform(0,9) in
        x5 < 0.5
    \end{lstlisting}
    \caption{Conditional Independent (Figure~\ref{fig:cond-benchmarks-a}): Program of size 5. Each variable depends on its immediate predecessor variable via a conditional branch. Each uniform distribution is distinct.}
\end{figure}

\begin{figure}
    \begin{lstlisting}
        let x1 = uniform(0,1) in
        let x2 = if x1 < 0.18 then uniform(0,2) else uniform(0,3) in
        let x3 = if x1 < 0.63 then uniform(0,4) else uniform(0,5) in
        let x4 = if x2 < 0.09 then uniform(0,6) else uniform(0,7) in
        let x5 = if x1 < 0.14 then uniform(0,8) else uniform(0,9) in
        x5 < 0.5
    \end{lstlisting}
    \caption{Conditional Random 1 (Figure~\ref{fig:cond-benchmarks-b}): Program of size 5. EAch variable depends on a randomly selected predecessor via a conditional branch. Note that this structure may lead to unused variables; that is, some previously defined $x_i$ may not influence any subsequent computation. Each uniform distribution is distinct.}
\end{figure}

\begin{figure}
    \begin{lstlisting}
        let x1 = uniform(0, 10) in
        let x2 = if x1 < 0.06 then x1 else x1 in
        let x3 = if uniform(0, 4) < 0.64 then x2 else x1 in
        let x4 = if x3 < 0.71 then uniform(0, 5) else x1 in
        let x5 = if uniform(0, 4) < 0.4 then x1 else uniform(0, 2) in
        x4 < 0.5
    \end{lstlisting}
    \caption{Conditional Random 2 (Figure~\ref{fig:cond-benchmarks-c}): Program of size 5. Each variable uses a randomized conditional expression where the guard and both branches may involve either a randomly selected predecessor variable or a fresh uniform distribution. The second parameter of the uniform distribution is a random integer between 0 and 10, so this may result in repeated bounds across different variables, meaning not all distributions are guaranteed to be unique.}
\end{figure}

\subsection{Alternating Guard Programs}
\begin{figure}
    \begin{lstlisting}
        let x1 = uniform(0,1) in
        let x2 = if x1 < 0.02 then uniform(0,2) else uniform(0,3) in
        let x3 = if x2 < 0.68 then uniform(0,4) else uniform(0,5) in
        let x4 = if x1 < 0.21 then uniform(0,6) else uniform(0,7) in
        let x5 = if x2 < 0.93 then uniform(0,8) else uniform(0,9) in
        x5 < 0.5
    \end{lstlisting}
    \caption{Alternating Guard 1 (Figure~\ref{fig:alt-benchmarks-a}): Program of size 5 with guard span 2. Guard variables in the comparisons cycle with a fixed span of 2. Note this may lead to unused variables, in particular those variables $x_i$ where $i$ is greater than the span.}
\end{figure}

\begin{figure}
    \begin{lstlisting}
        let x1 = uniform(0,1) in
        let x2 = if x1 < 0.38 then uniform(0,2) else x1 in
        let x3 = if x2 < 0.71 then uniform(0,4) else x2 in
        let x4 = if x1 < 0.42 then uniform(0,6) else x3 in
        let x5 = if x2 < 0.57 then uniform(0,8) else x4 in
        x5 < 0.5
    \end{lstlisting}
    \caption{Alternating Guard 2 (Figure~\ref{fig:alt-benchmarks-b}): Program of size 5 with guard span 2. Guard variables in the comparisons cycle with a fixed span of 2. Unlike Alternating Guard 1, all variables are used in subsequent branches, either as guards or else branch expressions. This construction ensures that every variable contributes to the final computation. Each uniform distribution generated is unique.}
\end{figure}

\begin{figure}
    \begin{lstlisting}
        let x1 = uniform(0,1) in
        let x2 = if x1 < 0.67 then uniform(0,2) else uniform(0,3) in
        let x3 = if x2 < 0.5 then x2 else uniform(0,4) in
        let x4 = if x1 < 0.5 then x3 else uniform(0,6) in
        let x5 = if x2 < 0.5 then x4 else uniform(0,8) in
        x5 < 0.5
    \end{lstlisting}
    \caption{Alternating Guard 3 (Figure~\ref{fig:alt-benchmarks-c}): Program of size 5 with guard span 2. Guard variables in the comparisons cycle with a fixed span of 2. Unlike Alternating Guard 1, all variables are used in subsequent branches, either as guards or if branch expressions. This construction ensures that every variable contributes to the final computation. Each uniform distribution generated is unique.}
\end{figure}

\begin{figure}
    \begin{lstlisting}
        let x1 = uniform(0,2) in
        let x2 = if x1 < 0.13 then uniform(0,2) else uniform(0,1) in
        let x3 = if x2 < 0.06 then uniform(0,8) else uniform(0,3) in
        let x4 = if x1 < 0.68 then uniform(0,4) else uniform(0,8) in
        let x5 = if x2 < 0.51 then uniform(0,3) else uniform(0,7) in
        x5 < 0.5
    \end{lstlisting}
    \caption{Random Alternating (Figure~\ref{fig:alt-benchmarks-d}): Program of size 5 with guard span 2. Guard variables in the comparisons cycle with a fixed span of 2. The second parameter of the uniform distribution is a random integer between 0 and 10, so this may result in repeated bounds across different variables, meaning not all distributions are guaranteed to be unique.}
\end{figure}

